\chapter{The five-slot lens}

\section{Building the lens}

The previous three chapters have been about reading. We toured five
artificial-life systems; we sat with the philosophical question of what
life is; we surveyed the contested vocabulary of emergence and novelty.
Now we build the lens.

The lens is a five-component decomposition of any artificial-life system.
Stated in plain English first:

\begin{enumerate}[leftmargin=2em,itemsep=2pt]
  \item \textbf{Substrate}: the material being simulated and its native
    dynamics.
  \item \textbf{Variation}: the rule that proposes changes to the
    entities living on the substrate.
  \item \textbf{Viability}: the rule that decides which proposed
    entities persist.
  \item \textbf{Interaction topology}: the structure of who-affects-whom.
  \item \textbf{Observer}: the measurement apparatus that scores what we
    can actually see.
\end{enumerate}

The names are abbreviated to $X$, $V$, $F$, $T$, $O$ in the paper. The
choice of letter is conventional and not worth defending. What matters is
that each system can be written down as a five-tuple
\[
  \mathcal{S} \;=\; \big(\mathbf{X},\, V,\, F,\, T,\, O\big),
\]
and that two systems described this way can be compared slot by slot.

This chapter introduces the slots in their typed form, defines the joint
state space on which the system iterates, writes out the tick update
rule, and re-reads the five systems from Chapter~1 through the lens.

\section{The substrate}

The substrate carries everything. It is the simulation material plus its
native update rule. Formally, a substrate is a triple
\[
  \mathbf{X} \;=\; (X,\, \Phi,\, \mu_0),
\]
where:

\begin{itemize}[leftmargin=1.5em,itemsep=2pt]
  \item $X$ is a \emph{measurable space} --- the state space, in plain
    English. For the Game of Life, $X = \{0,1\}^{H\times W}$, a grid of
    binary cells. For Lenia, $X = [0,1]^{H\times W}$, a grid of unit-
    interval values. For Boids with $N$ particles in $\mathbb{R}^d$,
    $X = (\mathbb{R}^d \times \mathbb{R}^d)^N$ (position-and-velocity
    pairs). For Tierra, $X$ is the contents of a circular memory of
    fixed size.
  \item $\Phi : X \to X$ is the \emph{substrate dynamics}: the function
    that updates the state by one tick. For the Game of Life, $\Phi$
    applies the birth-and-survival rule simultaneously across the grid.
    For Lenia, $\Phi$ convolves the field with the kernel, applies the
    growth function, and clamps to $[0,1]$. For Boids, $\Phi$ updates
    each particle's velocity by the steering rules and its position by
    the new velocity. If the dynamics is stochastic, we let $\Phi$ map
    to a probability distribution over $X$ instead of a single $X$, with
    the rest of the apparatus unchanged.
  \item $\mu_0$ is the \emph{initial measure} --- a non-negative measure
    on $X$, usually concentrating on a single state for one-shot runs,
    or distributed over starting states for ensemble runs. Treat $\mu_0$
    as ``the initial population.''
\end{itemize}

That is the substrate. Nothing about who interacts with whom, nothing
about variation or selection, nothing about what an observer sees. The
substrate is just the simulation kernel and its natural update.

If the words \emph{measurable space} or \emph{measure} are new to you,
flip to the math primer at the back of the book. The short version: a
measurable space is a state space with a notion of which subsets have a
well-defined size, and a measure is a way of assigning sizes to those
subsets. Probability measures sum to one; positive measures need not.
You will mostly think of them as ``distributions over states.''

\section{Variation}

Variation is what generates new candidates. In the language of the lens,
\[
  V \;:\; X \times \Theta \;\to\; \mathcal{P}(X),
\]
i.e.\ $V$ takes a state $x \in X$ and a parameter $\theta \in \Theta$ and
returns a probability distribution over next states. The parameter space
$\Theta$ holds the knobs --- mutation rate, recombination probability,
gradient step size, LLM prompt --- and the codomain is a distribution
over $X$ so that stochastic, deterministic, and gradient-based variation
all fit the same signature.

Examples:

\begin{itemize}[leftmargin=1.5em,itemsep=2pt]
  \item \textbf{Gaussian mutation}: $V(x, \sigma) = \mathcal{N}(x,
    \sigma^2 I)$. Add Gaussian noise.
  \item \textbf{Bit-flip}: each bit of $x$ flips with probability $p$.
  \item \textbf{Recombination}: pick two parents, output a child whose
    bits or fields are drawn from either parent.
  \item \textbf{Gradient descent}: $V(x, \eta) = \delta_{x - \eta \nabla
    L(x)}$, a Dirac mass on a single next state. Deterministic variation
    is a special case of the general signature.
  \item \textbf{LLM-prompt mutation}: $V(x, \theta)$ samples a new
    sequence by querying a language model with $x$ in context, biased by
    $\theta$.
\end{itemize}

The signature unifies these. A variation operator is a stochastic
function from current state and parameters to a distribution over next
states. We will see in Chapter~5 that variation also has a
\emph{composition operator} $\diamond_\rho$ for coupling two variation
channels (say, genetic and cultural).

\section{Viability}

Viability filters the population. After variation has proposed candidate
entities, viability decides who stays. Formally,
\[
  F \;:\; \mathcal{M}_+(X) \;\to\; \mathcal{M}_+(X),
\]
i.e.\ $F$ takes a positive measure on $X$ (a weighted population of
entities) and returns a positive measure on $X$ (a reweighted or
sub-sampled version). The codomain is $\mathcal{M}_+(X)$ rather than
$\mathcal{P}(X)$ because $F$ may shrink the total mass: extinction is
allowed, and post-filter populations need not be probability
distributions.

The interesting fact, and the one that will occupy Chapter~6, is that
there are several substantively different ways to write down $F$ in the
literature:

\begin{itemize}[leftmargin=1.5em,itemsep=2pt]
  \item \textbf{Markov-blanket viability}: an entity is viable iff it
    admits a sparse-coupling decomposition into internal and external
    states.
  \item \textbf{Autopoietic closure}: an entity is viable iff its
    supporting set is closed under its own reaction network.
  \item \textbf{RAF set viability}: an entity is viable iff its
    reaction set is reflexively autocatalytic and food-generated.
  \item \textbf{Minimal Criterion Coevolution}: an entity is viable iff
    it satisfies a minimal qualifying predicate.
\end{itemize}

These are not the same. The book treats them as four distinct
formalisms and does not assume they collapse to a common closure (that
would be Conjecture~C1, see Chapter~9 and \texttt{RESULTS.md}). For
purposes of the lens, the slot is filled by one of them, named
explicitly.

\section{Interaction topology}

Interaction topology is the rule of who-affects-whom. It is \emph{not}
the geometry of the substrate. The substrate may live on a square
lattice; the interaction topology may be a different graph entirely, say
the graph of gliders that have collided, or the bipartite graph of
language-model agents talking to Lenia creatures.

Formally, the interaction topology is a family $T = (T_t)_{t \geq 0}$,
where each $T_t$ is a (hyper)graph on the entity set at time $t$, with
weighted edges in some semiring (typically $\mathbb{R}_{\geq 0}$).
Concrete instances:

\begin{itemize}[leftmargin=1.5em,itemsep=2pt]
  \item Moore neighbourhood on a lattice.
  \item $k$-nearest-neighbour graph on $\mathbb{R}^d$ particles.
  \item Static social graph.
  \item Random donor--recipient pairing per tick.
  \item Stigmergic indirect coupling, where interactions pass through
    the substrate rather than directly between entities (an ant trail).
\end{itemize}

When several coupling regimes coexist (a chemical-signalling layer and
a pheromone-trail layer), $T_t$ becomes a \emph{multiplex}: a disjoint
union of layer-specific (hyper)graphs with weights $(\alpha_\ell)$
summing to one. The mass-conservation requirement is what makes the
mathematics work; we will revisit it in Chapter~7.

\section{The observer}

The observer is the measurement apparatus. It is not another driver of
the dynamics. It watches trajectories and reports what it sees.

In the stateless case,
\[
  O \;:\; Z^w \;\to\; \mathbb{R}^k,
\]
where $Z = X \times \mathcal{M}_+(X)$ is the joint state space (defined
below) and $w \in \mathbb{N} \cup \{\infty\}$ is a window length: how
many ticks of history the observer consumes per evaluation. The output
$\mathbb{R}^k$ holds the observed quantities --- novelty, complexity,
learnability, value, or whatever the observer measures.

Many useful observers are not stateless. A novelty-search observer
maintains an archive; a Quality-Diversity observer maintains a behaviour
grid; a foundation-model observer carries a learned embedding. For
these, we allow an observer state $a_t \in A_O$ and write
\[
  O_s \;:\; Z^w \times A_O \to \mathbb{R}^k \times A_O,
\]
i.e.\ the observer takes a window of trajectory and its current state,
returns a score vector and an updated state.

Three observer families recur (full treatment in Chapter~7):

\begin{itemize}[leftmargin=1.5em,itemsep=2pt]
  \item \textbf{Density-distance observers}, e.g.\ novelty search and
    QD.
  \item \textbf{Foundation-model embedding observers}, e.g.\ ASAL,
    OMNI-EPIC.
  \item \textbf{Information-theoretic observers}, e.g.\ mutual
    information, transfer entropy, effective information.
\end{itemize}

The observer slot is the most consequential one for novelty: novelty is
always relative to the observer. The same trajectory can be highly novel
under one observer and trivially repetitive under another.

\section{The joint state space and the iteration rule}

We have five typed objects. Now we wire them together.

The system iterates on the \emph{joint state space}
\[
  Z \;:=\; X \times \mathcal{M}_+(X).
\]
An element $(x, \mu) \in Z$ is a pair: a substrate state $x$ and a
population measure $\mu$ over the entities currently living on the
substrate. For substrate-as-population paradigms like Lenia or NCA,
$\mu$ is implicit in $x$ and the second coordinate is ignored. For
population-of-particles or population-of-programs paradigms, $\mu$ is
genuine.

Each slot induces a lifted operator on $Z$:
\begin{align*}
  \widetilde\Phi(x,\mu)   &= \big(\Phi(x),\, \Phi_*\mu\big), \\
  \widetilde V_T(x,\mu)   &= \big(x,\, V_T(\mu;\, x)\big), \\
  \widetilde S_F(x,\mu)   &= \big(x,\, F(\mu)\big),
\end{align*}
where $\Phi_*\mu$ is the push-forward of $\mu$ along $\Phi$ (informally:
each entity moves wherever $\Phi$ sends it), and $V_T(\mu; x)$ is the
population-level variation kernel produced by $V$ and $T$ together
(Chapter~5). The tick rule is the composition:
\[
  (x_{t+1},\, \mu_{t+1}) \;=\;
  \big(\widetilde S_F \circ \widetilde V_T \circ \widetilde\Phi\big)
  (x_t,\, \mu_t).
\]
In words: the substrate dynamics fires, then variation proposes new
candidates mediated by topology, then viability filters the result.

The observer is not part of the tick rule. It runs alongside, consuming
a window of trajectory at each tick:
\[
  (o_t,\, a_{t+1}) \;=\; O_s\big((x,\mu)_{t-w:t+1},\, a_t\big).
\]
The observer does not, by default, feed back into the iteration. Some
systems do feed observer output back (novelty search uses observer
state to shape variation), and the lens accommodates this through the
$\rho$-coupling machinery of Chapter~5; the default, however, is that
the observer is read-only.

\subsection*{Iteration order is data, not behaviour}

The order $\widetilde S_F \circ \widetilde V_T \circ \widetilde\Phi$ is
not the only sensible choice. Some systems apply selection before
variation, some run variation across multiple sub-ticks, some interleave
the operators with the observer. The lens does not impose a single
order. It does require you to \emph{name} the order you are using:
\[
  \text{iteration\_order}
  = (\text{``dynamics''},\, \text{``variation''},\, \text{``viability''})
\]
or some permutation thereof. The default above is the one most commonly
used. A system that departs from it should declare the deviation.

\section{Re-reading the five systems}

Chapter~1 introduced five systems. The lens now lets us write them down
slot by slot.

\subsection*{Game of Life}

\begin{itemize}[leftmargin=1.5em,itemsep=2pt]
  \item $\mathbf{X}$: $X = \{0,1\}^{H\times W}$; $\Phi$ is the GoL
    update; $\mu_0$ concentrates on a chosen initial pattern.
  \item $V$: trivial (no variation). The Game of Life is a pure
    dynamics.
  \item $F$: trivial (no filtering --- every state proceeds).
  \item $T$: the Moore neighbourhood on $L = \mathbb{Z}/H \times
    \mathbb{Z}/W$. Interaction topology coincides with substrate
    geometry here.
  \item $O$: whatever you want. Counting gliders. Counting live cells.
    Checking for stable configurations. The lens does not specify.
\end{itemize}

The Game of Life is a stripped-down case: substrate-only. That is why
it is useful pedagogically; everything except $\mathbf{X}$ is empty.

\subsection*{Boids}

\begin{itemize}[leftmargin=1.5em,itemsep=2pt]
  \item $\mathbf{X}$: $X = (\mathbb{R}^d \times \mathbb{R}^d)^N$
    (position + velocity per boid); $\Phi$ applies alignment, cohesion,
    separation; $\mu_0$ is random initial positions and velocities.
  \item $V$: trivial.
  \item $F$: trivial.
  \item $T$: $k$-nearest neighbours, or $\varepsilon$-ball, on
    positions. Substrate is $\mathbb{R}^d$; interaction topology is
    decidedly not the manifold structure of $\mathbb{R}^d$ --- it is a
    graph on particles.
  \item $O$: a flock-shape measure, a cohesion measure, or nothing.
\end{itemize}

Boids highlights the topology slot: changing the neighbourhood radius
changes flocking behaviour without touching the substrate or the
dynamics.

\subsection*{Lenia}

\begin{itemize}[leftmargin=1.5em,itemsep=2pt]
  \item $\mathbf{X}$: $X = [0,1]^L$ for $L$ a 2D torus; $\Phi$ is
    kernel-convolution-and-growth; $\mu_0$ supplies an initial field.
  \item $V$: a parameter mutation, applied to kernel weights when
    searching the parameter space.
  \item $F$: an entity-level filter on Lenia ``creatures'' (e.g.\
    autopoietic-closure on the local reaction network).
  \item $T$: the lattice topology by default; can be enriched with
    entity-spatial-overlap graphs in coupled systems.
  \item $O$: many choices. CLIP-embedding novelty (ASAL); the $\Omega$
    metric on attractors; the number of creatures.
\end{itemize}

Lenia is the book's running example because it has non-trivial content
in every slot.

\subsection*{Tierra}

\begin{itemize}[leftmargin=1.5em,itemsep=2pt]
  \item $\mathbf{X}$: $X$ = the contents of a circular memory of fixed
    size; $\Phi$ = the virtual-machine instruction cycle.
  \item $V$: bit-flips during replication.
  \item $F$: ``has the program reproduced?'' --- a minimal-criterion
    predicate, with memory contention culling the failures.
  \item $T$: implicit, via memory adjacency.
  \item $O$: distinct lineage count, ecology composition over time.
\end{itemize}

Tierra demonstrates the viability slot most clearly: a minimal
criterion (``replicates'') is enough to drive an entire ecology when
combined with variation and topology.

\subsection*{Evolutionary computation (GA)}

\begin{itemize}[leftmargin=1.5em,itemsep=2pt]
  \item $\mathbf{X}$: $X$ = the genome space; $\Phi$ = identity (genomes
    do not update on their own).
  \item $V$: crossover + mutation.
  \item $F$: fitness-proportional selection.
  \item $T$: complete graph (any genome can mate with any other), or a
    structured topology in spatially-embedded GAs.
  \item $O$: best-fitness-so-far, or a coverage measure on the
    behaviour space.
\end{itemize}

The genetic algorithm is the textbook case of \emph{variation +
viability are everything}: the substrate has trivial dynamics, and the
work is done by $V$ and $F$.

\section{Entities, briefly}

We have used the word \emph{entity} in passing. Now we can be precise
about it.

The substrate carries a state. The system's update rule does not, by
default, name pieces of the state as ``things.'' Calling a region of the
Lenia field a ``creature'' or a clump of boids a ``flock'' is a choice
made on top of the substrate state, not something the substrate hands
you.

The lens makes the choice explicit. An \emph{entity} is a pair
\[
  e \;=\; (S,\, \pi),
\]
where $S$ is a \emph{supporting set} (the substrate sites that the
entity occupies) and $\pi$ is a \emph{scale projection} (the description
we are committing to). For a Lenia creature, $S$ is the set of grid
sites above a mass threshold and $\pi$ is, say, the centroid-and-mass
map. For a Boids flock, $S$ is the set of particles in a cluster and
$\pi$ is the centroid and the inertia tensor. For an Avida program, $S$
is the memory region holding the program and $\pi$ is the identity
(the raw instruction sequence).

The supporting set $S$ tells you \emph{what the entity is made of}. The
scale projection $\pi$ tells you \emph{at what level of description we
are reading it}. Two readers can agree on $S$ and disagree on $\pi$; the
lens makes that disagreement visible.

There is an additional condition --- informational closure --- that an
$(S, \pi)$ pair should satisfy to count as an entity: most of the
information predictive of $\pi$'s future should be internal to $S$,
conditional on the outside. The full statement is in the paper and uses
mutual information. The intuition is the one we built up in Chapter~2:
the inside of a thing is informationally distinguishable from its
outside.

\section{What the lens is and is not}

The lens is a vocabulary. It is a place to put each modelling choice so
that two systems built by two people can be compared. It does not
arbitrate between candidate values of any slot, and it does not impose
a particular philosophy of life or emergence.

Three properties to internalise before we go further:

\begin{itemize}[leftmargin=1.5em,itemsep=2pt]
  \item \textbf{The slots are typed.} Each slot has a signature, and
    valid systems are those whose components match the signatures.
    Swapping the type of a slot (e.g.\ writing a viability filter as
    $X \to X$ rather than $\mathcal{M}_+(X) \to \mathcal{M}_+(X)$) is a
    category error, not just a stylistic difference.
  \item \textbf{The slots are independent.} You can change any one slot
    without rewriting the others. Variation does not know about
    topology; topology does not know about viability. The interfaces
    are minimal.
  \item \textbf{The slots compose.} Substrates compose via
    $\boxtimes_\kappa$ (Chapter~5); variation channels compose via
    $\diamond_\rho$ (Chapter~5); topologies compose into multiplexes
    (Chapter~7). Composition lets you build hybrid systems without
    abandoning the lens.
\end{itemize}

The next four chapters develop each slot in detail, with worked
examples and the formal machinery the slot needs. Chapter~5 is
substrate and variation. Chapter~6 is viability with its four
formalisms. Chapter~7 is topology and observers. Chapter~8 brings them
all together for the worked composite example and the quantitative
measurement of emergence and novelty.

\section*{To think about}

\begin{enumerate}
  \item Pick a system you wrote recently (anything with rules: a board
    game, a recommendation system, a physics simulation). Write it down
    as a five-tuple. Which slot was hardest to identify?
  \item The lens treats the observer as separate from the iteration
    rule. Why is this the right design choice? What would go wrong if
    we folded the observer into the tick update?
  \item For the Game of Life, $V$ and $F$ are trivial. What change to
    the rule would make a non-trivial $F$ natural? A non-trivial $V$?
  \item Two genetic algorithms with the same fitness function but
    different mutation rates are distinguished by which slot? Two GAs
    with the same operators but different population topologies?
  \item Suppose you stripped Boids of its three rules and kept only the
    interaction topology and the substrate. What kind of system would
    you have? Is it interesting?
\end{enumerate}
